\FloatBarrier
\section{本章小结}
\label{sec:integral-theorems-summary}

散度与旋度是边界累计的局部密度。连续可微向量场的散度$\sum_j\partial_jF_j$等于收缩立方体或小球上单位体积的外向净通量，也是Jacobian的迹；正值表示源，负值表示汇，散度为零只表示流入与流出平衡，不表示场为零。

三维场的旋度沿单位向量$\bm n$的分量，等于垂直于$\bm n$的微小圆周上逆时针环流除以面积的极限；平面场对应的标量是$Q_x-P_y$。

旋度与流线是否弯曲是两回事：剪切流的流线是直线而旋度不为零，涡旋场在原点以外流线弯曲而旋度为零。

Laplace算子$\Delta v=\operatorname{div}\nabla v$是Hessian的迹，度量一点的值与周围平均值之差；$C^2$时$\operatorname{curl}\nabla v=\bm0$与$\operatorname{div}\operatorname{curl}\bm F=0$来自混合偏导数的交换。

三个积分公式都由同一机制得到：在小块上建立局部关系，相容方向使内部公共边界的贡献相消。Green公式把平面区域上$Q_x-P_y$的积分化为正向边界上的环流，它的通量形式把散度的积分化为外向通量；正向边界要求区域总在行进方向的左侧，孔洞边界因而是顺时针。

Gauss公式把空间区域上散度的积分化为闭曲面上的外向通量。Stokes公式把定向曲面上旋度的通量化为诱导方向下的边界环流，平面情形还原为Green公式；具有同一有向边界的两张曲面给出相同的旋度通量，前提是场在两张曲面附近都连续可微。

区分闭曲面与有边界曲面：前者配外向通量和散度，后者配边界环流和旋度。由乘积法则与Gauss公式得到多元分部积分$\int_\Omega u\Delta v=\int_{\partial\Omega}u\nabla v^\top\bm n-\int_\Omega\nabla u^\top\nabla v$，交换$u,v$相减得到Green第二恒等式。边界项是否消失取决于边界条件：边界值为零时它消失，并由此推出Laplace方程的解由边界值唯一确定。

使用这些公式时要逐项检查：场在整个积分对象的闭包附近连续可微，边界分片规则，法向与边界方向相容。穿孔平面上的涡旋场与原点处的点源场说明，局部的零旋度或零散度不能越过场的奇点推出零环流或零通量。放宽光滑性或边界规则性，需要另行建立弱导数与相应的积分公式。
